% Mathematics & Research Notes - LaTeX template
% 编译方式：XeLaTeX（Overleaf 中 Menu -> Compiler -> XeLaTeX）
\documentclass[12pt,a4paper]{ctexart}

% 页面、数学环境和链接
\usepackage[margin=2.5cm]{geometry}
\usepackage{amsmath,amssymb,amsthm,mathtools}
\usepackage{graphicx}
\usepackage{booktabs}
\usepackage{hyperref}
\hypersetup{colorlinks=true,linkcolor=blue,citecolor=blue,urlcolor=blue}

% 定理环境
\theoremstyle{plain}
\newtheorem{theorem}{定理}[section]
\newtheorem{lemma}[theorem]{引理}
\newtheorem{proposition}[theorem]{命题}
\theoremstyle{definition}
\newtheorem{definition}[theorem]{定义}
\newtheorem{example}[theorem]{例}
\theoremstyle{remark}
\newtheorem{remark}[theorem]{注}
\renewcommand{\proofname}{证明}

% 常用记号
\newcommand{\R}{\mathbb{R}}
\newcommand{\dd}{\mathrm{d}}
\newcommand{\norm}[1]{\left\lVert #1\right\rVert}

\title{数学与科研笔记\\\large Mathematics \& Research Notes}
\author{作者姓名（请修改）}
\date{\today}

\begin{document}
\maketitle
\tableofcontents
\newpage

\section{研究背景与问题}
这里记录问题背景、文献来源、研究目标和基本假设。

\section{模型与基本定义}
考虑自治动力系统
\begin{equation}\label{eq:system}
    \dot{x}=f(x),\qquad x\in\R^n.
\end{equation}

\begin{definition}[平衡点]
若 $f(x^*)=0$，则称 $x^*$ 为系统 \eqref{eq:system} 的平衡点。
\end{definition}

\section{主要结果}
\begin{theorem}[线性化稳定性判据]\label{thm:stability}
设 $f$ 在平衡点 $x^*$ 的邻域内连续可微。若雅可比矩阵
$Df(x^*)$ 的全部特征值均具有严格负实部，则 $x^*$ 局部渐近稳定。
\end{theorem}

\begin{proof}
令 $y=x-x^*$，系统可写为
\begin{equation}
    \dot y=Ay+r(y),\qquad A=Df(x^*),\quad
    \lim_{\norm{y}\to0}\frac{\norm{r(y)}}{\norm{y}}=0.
\end{equation}
由线性化稳定性定理可得结论。实际写作时请在此补充完整证明或引用。
\end{proof}

\begin{remark}
如果存在特征值的实部等于零，上述判据不能直接决定非线性系统的稳定性。
\end{remark}

\section{示例：二维种群动力学}
一个简单的相互竞争模型是
\begin{equation}
\begin{cases}
\dot S=r_s S\left(1-\dfrac{S+\alpha R}{K}\right),\\[6pt]
\dot R=r_r R\left(1-\dfrac{R+\beta S}{K}\right).
\end{cases}
\end{equation}
这里 $S,R$ 表示两类种群数量，$r_s,r_r$ 为增长率，$K$ 为环境容纳量。

\section{数值模拟与图表}
\begin{table}[htbp]
\centering
\caption{参数表示例（数值仅供演示）}
\begin{tabular}{lll}
\toprule
参数 & 含义 & 示例值\\
\midrule
$r_s$ & 第一种群增长率 & $0.1$\\
$r_r$ & 第二种群增长率 & $0.08$\\
$K$ & 环境容纳量 & $1$\\
\bottomrule
\end{tabular}
\end{table}

% 如需插图，将图片放在 figures 文件夹下，取消下面各行的注释：
% \begin{figure}[htbp]
% \centering
% \includegraphics[width=0.75\textwidth]{figures/trajectory.png}
% \caption{轨迹示意图}
% \label{fig:trajectory}
% \end{figure}

\section{结论与待解决问题}
\begin{itemize}
    \item 本节总结已经得到的数学结论。
    \item 记录尚需验证的假设及数值问题。
    \item 记录下一步研究方向。
\end{itemize}

\begin{thebibliography}{9}
\bibitem{perko}
L. Perko, \textit{Differential Equations and Dynamical Systems}, 3rd ed., Springer, 2001.
\end{thebibliography}

\end{document}

